% Capítulo 3: Metodología
% Propósito: Detallar el enfoque metodológico, diseño de investigación,
% técnicas e instrumentos de recolección de información y la metodología
% de desarrollo de software seleccionada.

\section{Metodología}

La metodología del proyecto establece el procedimiento que se seguirá para obtener y analizar la información necesaria para el desarrollo de TravelDLS. En esta sección se define el enfoque y tipo de investigación, el procedimiento utilizado para seleccionar las investigaciones académicas que sustentan el proyecto, la población y muestra, las variables de estudio y las técnicas e instrumentos de recolección de datos. De esta manera, se busca mantener una relación clara entre el problema identificado, la información obtenida de los participantes y las decisiones que posteriormente se tomarán durante el diseño y desarrollo de la plataforma.


\subsection{Enfoque de la Investigación}

La investigación de TravelDLS se desarrollará bajo un \textbf{enfoque cuantitativo}, debido a que la información principal será obtenida mediante preguntas cerradas dirigidas a administradores o propietarios y conductores vinculados con empresas de transporte terrestre de carga pesada. Los resultados serán organizados mediante frecuencias, porcentajes, tablas y gráficos con el propósito de identificar prácticas actuales, necesidades y problemas relacionados con la gestión de viajes, seguimiento, comunicación y planificación de rutas.

\textcite{queiros2017strengths} señalan que la investigación cuantitativa busca obtener mediciones precisas y confiables que permitan posteriormente realizar un análisis estadístico de los datos. Este enfoque resulta adecuado para TravelDLS porque permitirá representar de forma numérica aspectos como las herramientas utilizadas por las empresas, las formas de seguimiento de los viajes, los medios empleados para planificar rutas y la frecuencia con que se presentan determinados problemas.

El cuestionario también incluirá algunas preguntas abiertas para permitir que los participantes expresen situaciones o dificultades que no hayan sido contempladas previamente dentro de las opciones de respuesta. Estas respuestas serán utilizadas como información complementaria para interpretar los resultados cuantitativos, pero no constituirán por sí mismas un enfoque cualitativo independiente, debido a que el análisis principal del estudio estará centrado en los datos cuantificables obtenidos mediante el instrumento.


\subsection{Tipo de Estudio}

Por su finalidad, TravelDLS corresponde a una \textbf{investigación aplicada}, ya que el estudio no se limita a describir una situación, sino que busca utilizar el conocimiento obtenido para desarrollar una solución dirigida a una necesidad concreta. \textcite{gulbrandsen2010concepts} explican que la investigación aplicada se encuentra orientada hacia objetivos prácticos específicos y a la utilización del conocimiento para atender problemas o necesidades determinadas. En TravelDLS, la información obtenida mediante la revisión académica y el estudio de campo será utilizada para establecer requerimientos y orientar el diseño y construcción de la plataforma.

Durante la etapa de diagnóstico, la investigación tendrá además un \textbf{alcance descriptivo}, debido a que se pretende conocer cómo se realizan actualmente determinadas actividades relacionadas con las operaciones de transporte de carga pesada. \textcite{siedlecki2020understanding} señala que los estudios descriptivos se utilizan para observar y describir las características de una población o fenómeno tal como se presentan, sin manipular las condiciones estudiadas. En este proyecto se describirán aspectos relacionados con la gestión de viajes, vehículos y conductores, los medios de comunicación utilizados, las formas de seguimiento y los procedimientos empleados para planificar las rutas.

El proyecto incorpora también un \textbf{componente de desarrollo tecnológico}, ya que los resultados de la investigación serán utilizados posteriormente para diseñar, construir, integrar y probar una plataforma web. \textcite{delacruz2016metodologia} plantea que la investigación tecnológica se encuentra orientada a la creación o mejora de procesos, sistemas y soluciones mediante la aplicación del conocimiento disponible. Esta característica se encuentra presente en TravelDLS porque el proceso investigativo servirá como base para desarrollar una solución tecnológica destinada a apoyar las operaciones de transporte terrestre de carga pesada.

De esta manera, el proyecto puede caracterizarse como una investigación \textbf{aplicada por su finalidad, descriptiva durante su etapa de diagnóstico y con un componente de desarrollo tecnológico} por el producto que se pretende construir.


\subsection{Procedimiento para la Selección de Investigaciones Relevantes}

Para seleccionar las investigaciones utilizadas como sustento del proyecto, se tomó como referencia el proceso de mapeo sistemático propuesto por \textcite{petersen2008systematic}. Esta metodología resulta adecuada para obtener una visión general de un campo de investigación, identificar los tipos de estudios existentes y reconocer áreas que han recibido mayor o menor atención.

En el caso de TravelDLS, el proceso comenzó con la definición de preguntas orientadas a temas relacionados con la gestión del transporte, trazabilidad, seguimiento, geolocalización, optimización de rutas y plataformas web aplicadas a procesos logísticos. Posteriormente, se realizaron búsquedas de publicaciones académicas utilizando términos relacionados con estos conceptos.

El procedimiento se apoya en las etapas establecidas por \textcite{petersen2008systematic}, quienes plantean como pasos principales la definición de las preguntas de investigación, la búsqueda de estudios primarios, el filtrado de los documentos encontrados, la clasificación mediante palabras clave y la extracción de información para construir el mapa del área estudiada (pp. 2--5).

Para determinar cuáles investigaciones podían aportar al proyecto se aplicaron criterios de inclusión y exclusión. Se consideraron principalmente artículos científicos, trabajos académicos y estudios relacionados con transporte, logística, gestión de flotas, seguimiento, geolocalización, planificación u optimización de rutas y desarrollo de plataformas tecnológicas relacionadas con estos procesos. Se descartaron documentos que solamente mencionaban estos conceptos de manera superficial, publicaciones duplicadas y materiales que no aportaban información relacionada con los objetivos del proyecto.

Este proceso también tomó en cuenta la recomendación de revisar otras secciones del documento cuando el resumen no proporcionara información suficiente para determinar su pertinencia. \textcite{petersen2008systematic} señalan que, cuando el abstract no permite identificar adecuadamente la contribución de una investigación, pueden revisarse apartados como la introducción o las conclusiones antes de realizar su clasificación.

A partir de los estudios seleccionados se identificaron palabras clave y conceptos recurrentes, los cuales permitieron organizar las investigaciones en categorías relacionadas con gestión logística, seguimiento y geolocalización, optimización de rutas, sistemas de gestión de transporte, plataformas SaaS y tecnologías utilizadas para el desarrollo de sistemas web.


\subsection{Población y Muestra}

La población del estudio estará conformada por personas vinculadas directamente con empresas dedicadas al transporte terrestre de carga pesada dentro de \textbf{[indicar el área geográfica definitiva del estudio]}. Para la recolección de información se considerarán principalmente dos grupos: administradores o propietarios de empresas transportistas y conductores que participan en la ejecución de los viajes.

La selección de estos participantes responde a que ambos grupos intervienen directamente en procesos que TravelDLS pretende apoyar. Los administradores o propietarios participan en actividades como la organización de viajes, asignación de vehículos y conductores, gestión de información y comunicación con los clientes. Por su parte, los conductores poseen experiencia directa sobre la recepción de información del viaje, utilización de rutas, herramientas de navegación, cambios durante el recorrido y comunicación con la empresa.

\textcite{taherdoost2016sampling} señala que el proceso de muestreo debe comenzar con la definición de la población objetivo y continuar con la identificación del marco muestral, la elección del método de muestreo y la determinación del tamaño de la muestra. Por esta razón, el número definitivo de participantes de TravelDLS será establecido una vez que se determine la población accesible dentro del territorio seleccionado.

En caso de que no se disponga de un registro completo de todos los administradores, propietarios y conductores que permita efectuar una selección aleatoria, se utilizará un \textbf{muestreo no probabilístico intencional}. Este procedimiento permitirá seleccionar participantes que cumplan con los perfiles establecidos para la investigación y que posean experiencia directa en las operaciones de transporte terrestre de carga pesada.

\textcite{taherdoost2016sampling} diferencia los métodos probabilísticos de los no probabilísticos y explica que la selección de la técnica debe considerar las características de la población, la disponibilidad de un marco muestral y las condiciones bajo las cuales se realiza la investigación. En TravelDLS, la elección definitiva del procedimiento de muestreo dependerá de la información disponible sobre las empresas y participantes que conformen la población de estudio.


\subsection{Operacionalización de Variables}

Para orientar la recolección y posterior análisis de la información se establecieron las principales variables de estudio relacionadas con TravelDLS. Estas variables permiten organizar los aspectos que serán observados durante el diagnóstico y posteriormente evaluados durante el desarrollo y las pruebas de la plataforma.

En este sentido, se consideran el sistema web TravelDLS, la gestión de operaciones de transporte, la localización y trazabilidad, la optimización de rutas y la transparencia para el cliente. Cada una de estas variables se divide en dimensiones e indicadores que permiten establecer de manera más concreta qué información será recopilada o evaluada durante el proyecto.

La relación entre estos elementos se presenta en la siguiente tabla:

\begin{table}[H]
	\centering
	\caption{Operacionalización de Variables}
	\label{tab:operacionalizacion}
	\small
	\renewcommand{\arraystretch}{1.3}
	\begin{tabularx}{\textwidth}{|>{\RaggedRight\arraybackslash}p{2.5cm}|>{\RaggedRight\arraybackslash}p{3.5cm}|>{\RaggedRight\arraybackslash}p{2.8cm}|>{\RaggedRight\arraybackslash}X|}
		\hline
		\textbf{Variable} & \textbf{Definición conceptual} & \textbf{Dimensiones} & \textbf{Indicadores} \\
		\hline
		
		\textbf{Sistema web TravelDLS} &
		Plataforma web SaaS destinada a centralizar y gestionar información relacionada con las operaciones de transporte de carga pesada. &
		Gestión de información \newline
		Gestión de operaciones \newline
		Integración &
		Registro y consulta de información \newline
		Viajes, vehículos, conductores y cargas gestionados \newline
		Módulos integrados y funcionamiento conjunto \\
		\hline
		
		\textbf{Gestión de operaciones de transporte} &
		Organización de la información y actividades necesarias para coordinar un servicio de transporte. &
		Administración \newline
		Coordinación \newline
		Servicios &
		Viajes, cargas, vehículos y conductores gestionados \newline
		Asignación de vehículos y conductores \newline
		Servicios publicados y contratados \\
		\hline
		
		\textbf{Localización y trazabilidad} &
		Capacidad de conocer la ubicación y evolución del servicio durante el traslado. &
		Geolocalización \newline
		Seguimiento \newline
		Trazabilidad &
		Ubicación actual de la unidad \newline
		Estado actualizado del viaje \newline
		Historial del recorrido y estados \\
		\hline
		
		\textbf{Optimización de rutas} &
		Proceso de comparar alternativas para obtener recorridos considerando distancia, tiempo u otros costos. &
		Distancia \newline
		Tiempo \newline
		Comparación de métodos \newline
		Eficiencia &
		Distancia total de la ruta \newline
		Tiempo estimado del recorrido \newline
		Resultados de Vecino Más Cercano, 2-opt y OR-Tools \newline
		Diferencia entre resultados obtenidos \\
		\hline
		
		\textbf{Transparencia para el cliente} &
		Disponibilidad de información sobre el servicio contratado durante su desarrollo. &
		Estado \newline
		Ubicación \newline
		Información económica &
		Estado de la carga \newline
		Ubicación del viaje \newline
		Precio final mostrado antes de aceptar \\
		\hline
		
	\end{tabularx}
\end{table}


\subsection{Técnicas e Instrumentos de Recolección de Datos}

Para obtener información directamente de los participantes se utilizará la \textbf{encuesta como técnica de recolección de datos} y un \textbf{cuestionario estructurado como instrumento}. \textcite{kelley2003good} señalan que la investigación mediante encuestas requiere definir claramente la población objetivo, establecer el procedimiento de selección de participantes, diseñar adecuadamente el instrumento y determinar previamente la forma en que serán analizados los datos obtenidos.

El cuestionario de TravelDLS estará organizado de acuerdo con el rol desempeñado por cada participante. Para los administradores o propietarios se incluirán preguntas relacionadas con la cantidad de vehículos y conductores, la forma en que se almacena la información, la asignación de viajes, el seguimiento de las unidades, la planificación de rutas, la comunicación con los clientes y la disposición para utilizar una plataforma web.

Para los conductores se recopilará información relacionada con la forma en que reciben los datos de sus viajes, las herramientas utilizadas para navegación, la planificación de recorridos con varios destinos, los cambios de ruta, la comunicación durante el viaje, la disponibilidad de teléfonos inteligentes y el acceso a Internet.

El instrumento estará compuesto principalmente por preguntas cerradas de selección, frecuencia y escalas de respuesta, lo que facilitará la organización de los resultados mediante frecuencias y porcentajes. También se incorporarán preguntas abiertas para que los participantes puedan señalar problemas o situaciones adicionales que no hayan sido consideradas previamente.

\textcite{artino2014developing} proponen un proceso sistemático para el diseño de cuestionarios que incluye la revisión de literatura, la definición de los aspectos que se pretende medir, la elaboración de los ítems y su revisión antes de la aplicación definitiva. Estos elementos serán considerados durante la preparación del instrumento utilizado en TravelDLS.


\subsubsection{Validación del Instrumento}

Antes de aplicar el cuestionario de manera definitiva, se realizará una \textbf{validación de contenido mediante juicio de expertos}. El propósito será revisar si las preguntas son claras, pertinentes y coherentes con los objetivos y variables establecidas para el estudio.

\textcite{boateng2018best} señalan que la validación de contenido permite evaluar si los ítems representan adecuadamente el dominio que se pretende estudiar y consideran la evaluación por expertos como una de las actividades que pueden utilizarse durante esta etapa del desarrollo de un instrumento.

Los expertos revisarán aspectos como la claridad de la redacción, la relación de las preguntas con los objetivos de la investigación, la pertinencia de las opciones de respuesta y la presencia de preguntas ambiguas o repetitivas. Las observaciones obtenidas serán utilizadas para realizar las modificaciones necesarias antes de aplicar el cuestionario a la población seleccionada.

No se establece todavía un número definitivo de expertos, debido a que esta cantidad será determinada posteriormente de acuerdo con las orientaciones metodológicas establecidas para el proyecto.


\subsubsection{Prueba Piloto}

Después de incorporar las observaciones realizadas durante la validación del instrumento, se efectuará una \textbf{prueba piloto} con un grupo reducido de personas que posean características similares a las de la población objetivo.

La prueba tendrá como finalidad detectar dificultades de comprensión, preguntas que puedan generar confusión, problemas en las opciones de respuesta y conocer el tiempo aproximado necesario para completar el cuestionario. \textcite{artino2014developing} señalan la importancia de realizar pruebas previas durante el desarrollo de cuestionarios para identificar posibles dificultades antes de su aplicación definitiva.

Asimismo, \textcite{boateng2018best} incluyen la preprueba de las preguntas como una etapa del desarrollo de instrumentos, debido a que permite comprobar si los ítems y sus respuestas son comprensibles para los participantes antes de proceder con la aplicación de mayor alcance.

El número de participantes de la prueba piloto será establecido posteriormente de acuerdo con las características de la población y las recomendaciones del tutor metodológico.


\subsection{Procesamiento y Análisis de los Datos}

Una vez finalizada la aplicación del cuestionario, las respuestas cerradas serán organizadas y analizadas mediante \textbf{estadística descriptiva}. Se calcularán frecuencias y porcentajes para identificar las respuestas predominantes y la frecuencia con que se presentan determinadas situaciones dentro de las operaciones de transporte estudiadas.

Los resultados serán representados mediante tablas y gráficos para facilitar su interpretación. Este procedimiento se encuentra relacionado con el enfoque cuantitativo definido para la investigación, debido a que permitirá transformar las respuestas obtenidas en datos que puedan ser comparados y resumidos.

Las respuestas abiertas serán revisadas de manera complementaria con el propósito de identificar problemas, necesidades o comentarios que no se encuentren representados directamente dentro de las opciones cerradas del cuestionario. Estas respuestas no serán utilizadas para establecer una fase cualitativa independiente, sino como apoyo para interpretar y contextualizar los resultados cuantitativos.

La información obtenida mediante este proceso servirá como insumo para la identificación y priorización de requerimientos funcionales y no funcionales de TravelDLS.


\subsection{Metodología de Desarrollo}
